\section{Mellin moments, harmonic collisions, and polylogarithmic transforms}
\label{da:section}

The compact density has a useful arithmetic structure. This section
identifies it explicitly and shows how the same polynomials connect the
spectral formulas to Gamma ratios, polygamma differences, generalized
harmonic numbers, and multiple polylogarithms. The coefficient identities
and height-one reductions used in this bridge are classical. Their role
here is to give computable consequences of the logarithmic density, not
to assert a new historical origin for those identities.

\subsection{The beta generating function of the density}

Retain the polynomials \(P_r\) of \eqref{sp:coincident-polynomial} and put
\[
 y(t)=\log\frac{1-t}{t},\qquad 0<t<1.
\]
Since \(P_r(y)=-\pi^{-1}\Im(y-i\pi)^r\) for real \(y\), coefficient
expansion gives the entire generating identity
\begin{equation}\label{da:Pgen}
 \sum_{r\ge1}P_r(y)\frac{w^r}{r!}
       =e^{wy}\frac{\sin\pi w}{\pi}.
\end{equation}
For complex \(w\), the right side is the analytic continuation of this
polynomial coefficient identity; no imaginary-part operation is applied
to a complex parameter \(w\).

\begin{theorem}[Every Mellin moment of the logarithmic density]
\label{da:mellin}
For \(\Re s>0\), define the ordinary integral
\[
 M_r(s)=\int_0^1 t^{s-1}P_r(y(t))\dd t,\qquad r\ge1.
\]
For \(w\) in a neighborhood of zero,
\begin{equation}\label{da:mellin-gen}
 \boxed{\displaystyle
  \sum_{r\ge1}M_r(s)\frac{w^r}{r!}
     =\frac{w\Gamma(s-w)}{\Gamma(s+1)\Gamma(1-w)}.}
\end{equation}
More generally, the underlying integral of \eqref{da:Pgen} equals this
Gamma ratio whenever \(\Re(s-w)>0\) and \(\Re(1+w)>0\).
\end{theorem}

\begin{proof}
On a sufficiently small closed disk about \(w=0\), the integrand and all
fixed \(w\)-derivatives are bounded by
\[
 C\,t^{\Re s-1-\eta}(1-t)^{-\eta}
       (1+|\log t|^j+|\log(1-t)|^j),
\]
with \(0<\eta<\min(1,\Re s)\). This is integrable. Thus
coefficient extraction may be interchanged with the integral.
Euler's beta integral and the Gamma reflection relation then give
\begin{align*}
 \int_0^1t^{s-1}e^{wy(t)}\frac{\sin\pi w}{\pi}\dd t
 &=\frac{\sin\pi w}{\pi}
      \frac{\Gamma(s-w)\Gamma(1+w)}{\Gamma(s+1)}\\
 &=\frac{w\Gamma(s-w)}{\Gamma(s+1)\Gamma(1-w)}.
\end{align*}
This proves both the generating identity and its stated domain.
\end{proof}

Let \(Y_n\) denote the complete exponential Bell polynomials, normalized by
\[
 \exp\left(\sum_{j\ge1}a_j\frac{w^j}{j!}\right)
      =\sum_{n\ge0}Y_n(a_1,\ldots,a_n)\frac{w^n}{n!},
 \qquad Y_0=1.
\]
Taking the logarithm of the Gamma ratio in \eqref{da:mellin-gen} gives
the finite polygamma formula
\begin{equation}\label{da:polygamma-moment}
 M_r(s)=\frac r{s}Y_{r-1}(a_1(s),\ldots,a_{r-1}(s)),
 \quad
 a_j(s)=(-1)^j\{\psi^{(j-1)}(s)-\psi^{(j-1)}(1)\}.
\end{equation}
In particular,
\begin{align*}
 M_1(s)&=\frac1s,\\
 M_2(s)&=\frac2s\{\psi(1)-\psi(s)\},\\
 M_3(s)&=\frac3s\bigl\{(\psi(s)-\psi(1))^2+
                                   \psi'(s)-\psi'(1)\bigr\}.
\end{align*}
All \(s\)-derivatives of these identities are ordinary logarithmic Mellin
moments, justified by the same endpoint majorants.

\begin{corollary}[Stirling moments and exact orthogonality]
\label{da:stirling}
Let \(\stir nr\) be the unsigned Stirling number of the first kind, with
\(\stir nr=0\) when \(n<r\). For every \(k\ge0\), \(r\ge1\),
\begin{equation}\label{da:stirling-moment}
 \boxed{\displaystyle
  \int_0^1 t^kP_r\!\left(\log\frac{1-t}{t}\right)\dd t
    =(-1)^{r-1}\frac{r!}{(k+1)!}\stir{k+1}{r}.}
\end{equation}
Thus \(P_r(y(t))\) is orthogonal to every polynomial in \(t\) of degree
less than \(r-1\), with respect to ordinary Lebesgue integration.
\end{corollary}

\begin{proof}
At \(s=k+1\), the right side of \eqref{da:mellin-gen} is
\[
 \frac{w(1-w)_k}{(k+1)!}.
\]
Expanding this finite polynomial using
\(u(u+1)\cdots(u+k)=\sum_r\stir{k+1}{r}u^r\)
gives \eqref{da:stirling-moment} with its sign. The asserted
orthogonality follows from the vanishing Stirling numbers.
\end{proof}

This moment formula also makes the harmonic content explicit:
\[
 \frac{\stir{k+1}{r}}{k!}
 =e_{r-1}\left(1,\frac12,\ldots,\frac1k\right),
\]
where \(e_j\) is the elementary symmetric polynomial, and \(e_0=1\).
Equivalently, \eqref{da:polygamma-moment} at \(s=k+1\) has
\(a_j=-(j-1)!H_k^{(j)}\). These finite formulas make exact moment
calculation possible without numerical quadrature.

\subsection{A transform from ordinary to height-one multiple polylogarithms}

For \(|z|<1\) and \(s\in\mathbb C\), use
\[
 \Li_s(z)=\sum_{n\ge1}\frac{z^n}{n^s}.
\]
The height-one multiple polylogarithm needed here can be defined directly
by the locally normally convergent series
\begin{equation}\label{da:heightone-definition}
 \Li_{s+1,\one{r-1}}(z)
   =\sum_{n\ge1}\frac{z^n}{n^{s+1}}H_{n-1}(\one{r-1}),
\end{equation}
where
\[
 H_m(\one d)=
   \sum_{m\ge n_1>\cdots>n_d\ge1}\frac1{n_1\cdots n_d},
 \qquad H_m(\one0)=1.
\]
Thus the nesting convention is specified rather than inferred from
notation.

\begin{theorem}[The logarithmic-density polylogarithm transform]
\label{da:polylog}
For every \(r\ge1\), every \(s\in\mathbb C\), and \(|z|<1\),
\begin{equation}\label{da:polylog-transform}
 \boxed{\displaystyle
 \int_0^1 P_r\!\left(\log\frac{1-t}{t}\right)
                \frac{\Li_s(zt)}t\dd t
   =(-1)^{r-1}r!\,\Li_{s+1,\one{r-1}}(z).}
\end{equation}
The integral and series are locally normal in \(s\). Consequently every
order derivative satisfies
\begin{equation}\label{da:polylog-order-jets}
 \int_0^1P_r(y(t))\frac{\partial_s^j\Li_s(zt)}t\dd t
  =(-1)^{r-1}r!\,
       \partial_s^j\Li_{s+1,\one{r-1}}(z),\qquad j\ge0.
\end{equation}
For integer \(s\ge1\), the ordinary endpoint value at \(z=1\) is
\begin{equation}\label{da:mzv-transform}
 \int_0^1P_r(y(t))\frac{\Li_s(t)}t\dd t
   =(-1)^{r-1}r!\,\zeta(s+1,\one{r-1}).
\end{equation}
\end{theorem}

\begin{proof}
For \(|z|\le R<1\) and \(s\) in a compact set, the polylogarithm
series and any fixed order derivative are dominated by
\(\sum_n R^n n^A(1+\log^j n)\). Near zero,
\(\Li_s(zt)/t\) remains bounded uniformly on these sets; near one
the only additional factor is an integrable polynomial in
\(\log(1-t)\). Absolute termwise integration and
\eqref{da:stirling-moment} give
\[
 (-1)^{r-1}r!\sum_{n\ge1}
       \frac{z^n}{n^s n!}\stir nr.
\]
The finite identity
\(\stir nr/(n-1)!=H_{n-1}(\one{r-1})\) turns this into
\eqref{da:heightone-definition}. The same domination justifies the
order derivatives.

For real \(0\le z<1\) and integer \(s\ge1\),
\(|\Li_s(zt)|\le\Li_s(t)\). At \(s=1\) this grows at most as
\(|\log(1-t)|\), and at \(s>1\) it is bounded near one.
Dominated convergence gives the endpoint integral. The positive series
in \eqref{da:heightone-definition} converges at one because its first
index is at least two, proving \eqref{da:mzv-transform}.
\end{proof}

\begin{corollary}[An all-order ordinary zeta integral]
\label{da:zeta-integral}
For every integer \(r\ge1\),
\begin{equation}\label{da:zeta-integral-formula}
 \int_0^1P_r\!\left(\log\frac{1-t}{t}\right)
       \frac{-\log(1-t)}t\dd t
     =(-1)^{r-1}r!\,\zeta(r+1).
\end{equation}
For example,
\begin{equation}\label{da:pi-four}
 \boxed{\displaystyle
 \int_0^1\left(3\log^2\frac{1-t}{t}-\pi^2\right)
             \frac{-\log(1-t)}t\dd t=\frac{\pi^4}{15}.}
\end{equation}
\end{corollary}

\begin{proof}
The all-one iterated integral is
\(\Li_{\one r}(z)=(-\log(1-z))^r/r!\), by induction using its
derivative and its zero value at \(z=0\). Integration once against
\(dz/z\), followed by \(u=-\log(1-z)\), gives
\[
 \zeta(2,\one{r-1})
  =\frac1{r!}\int_0^1\frac{(-\log(1-z))^r}{z}\dd z
  =\frac1{r!}\int_0^\infty\frac{u^r}{e^u-1}\dd u
  =\zeta(r+1).
\]
The last equality follows by expanding \(1/(e^u-1)\) into its
positive exponential series and integrating. Apply
\eqref{da:mzv-transform} at \(s=1\). The case \(r=3\), together with
\(\zeta(4)=\pi^4/90\), gives \eqref{da:pi-four}.
\end{proof}

\subsection{Every confluent logarithmic coefficient}

For \(a>0\), \(\beta\ge0\), define the convergent resolvents
\[
 R_k(a;\beta)=\sum_{n\ge0}\frac1{((n+a)^2+\beta)^k}
   =\frac{(-1)^{k-1}}{(k-1)!}\partial_\beta^{\,k-1}Q_a(\beta),
 \qquad k\ge1.
\]
For \(r\) independent displacements satisfying
\(|\delta_j|<a^2+\beta\), the ordinary product sum has the convergent
collision expansion
\begin{align}
 &\sum_{n\ge0}\prod_{j=1}^r
     \log\left(1+\frac{\delta_j}{(n+a)^2+\beta}\right)\notag\\
 &\quad=
 \sum_{k_1,\ldots,k_r\ge1}
   \frac{(-1)^{k_1+\cdots+k_r-r}
             \delta_1^{k_1}\cdots\delta_r^{k_r}}
        {k_1\cdots k_r}\,
       R_{k_1+\cdots+k_r}(a;\beta).
 \label{da:independent-collision}
\end{align}
All branches are the ones obtained from the convergent logarithm series
at zero. On smaller closed polydisks, the series is absolutely and normally
convergent: the absolute logarithm series is bounded by a constant times
\(((n+a)^2+\beta)^{-1}\) in each factor, and the product is summable.
This proves the interchange and permits every fixed parameter derivative.

When all displacements equal \(\delta\), the classical logarithm-power
coefficients give a particularly short all-order formula:
\begin{equation}\label{da:confluent-stirling}
 \boxed{\displaystyle
 \sum_{n\ge0}\log^r\left(1+\frac{\delta}{(n+a)^2+\beta}\right)
   =r!\sum_{k\ge r}\frac{(-1)^{k-r}}{k!}\stir kr
                          \delta^kR_k(a;\beta).}
\end{equation}
For \(r=2\), this recovers the harmonic coefficient in the inherited
second-jet collision expansion \cite{ProveItHigherDefects}. For \(r=3\),
\begin{align}
 &\sum_{n\ge0}\log^3\left(1+\frac{\delta}{(n+a)^2+\beta}\right)\notag\\
 &\qquad=
 3\sum_{k\ge3}\frac{(-1)^{k-3}\delta^k}{k}
       \{H_{k-1}^2-H_{k-1}^{(2)}\}R_k(a;\beta).
 \label{da:cubic-harmonic}
\end{align}
At \(\beta=0\), \(R_k(a;0)=\zeta(2k,a)\); at an integer shift
\(a=M\), it is \(\zeta(2k)-H_{M-1}^{(2k)}\).
Thus the compact density, the convergent mass expansion, and the
polygamma expression for \(R_k\) are three independently useful
representations of the same all-order logarithmic sum.

\subsection{Half-lattices and bilateral polylogarithmic kernels}

Digamma reflection gives the particularly simple half-lattice kernel
\begin{equation}\label{da:half-lattice}
 Q_{1/2}(t)=\frac{\pi}{2\sqrt t}\tanh(\pi\sqrt t),
 \qquad Q_{1/2}(0)=\frac{\pi^2}{2}.
\end{equation}
For example the cubic repeated-mass identity becomes
\begin{align}
 &\sum_{n\ge0}\log^3
    \frac{(n+\frac12)^2+\nu}{(n+\frac12)^2+\beta}\notag\\
 &\quad=\frac{\pi}{2}\int_\beta^\nu
    \left\{3\log^2\frac{\nu-t}{t-\beta}-\pi^2\right\}
                  \frac{\tanh(\pi\sqrt t)}{\sqrt t}\dd t,
       \qquad \nu>\beta\ge0 .
 \label{da:half-cubic}
\end{align}
The apparent endpoint at \(t=0\) is removable in the hyperbolic kernel,
so this remains an ordinary integral when \(\beta=0\).

For \(0<a<1\), define the bilateral resolvent
\[
 Q_a^{\mathrm{bil}}(t)=Q_a(t)+Q_{1-a}(t).
\]
Reflection and the convergent geometric Fourier series give
\begin{align}
 Q_a^{\mathrm{bil}}(t)
 &=\frac{\pi}{\sqrt t}
     \frac{\sinh(2\pi\sqrt t)}
          {\cosh(2\pi\sqrt t)-\cos(2\pi a)}\notag\\
 &=\frac{\pi}{\sqrt t}
       \left\{1+2\Re\Li_0(e^{-2\pi\sqrt t+2\pi ia})\right\},
             \qquad t>0. \label{da:bilateral-kernel}
\end{align}
Here \(\Li_0(z)=z/(1-z)\), so the kernel is elementary.
The bilateral full spectral jet
\(J_{N,r}^{\mathrm{bil}}(a)=J_{N,r}(a)+J_{N,r}(1-a)\) therefore is
\begin{align}
 J_{N,r}^{\mathrm{bil}}(a;\boldsymbol\mu)
 ={}&(2N)^r\{\zeta^{(r)}(0,a)+\zeta^{(r)}(0,1-a)\}\notag\\
 &+(-1)^r\int_0^\infty
          \rho_{N,r}(t)Q_a^{\mathrm{bil}}(t)\dd t.
 \label{da:bilateral-jet}
\end{align}
This is a finite polylogarithmic-integrand description at every order,
while retaining the displayed mass-zero Hurwitz derivatives. It does not
assert a finite evaluation of the remaining mass integral.

There is also an explicit anchored shift primitive of the bilateral
kernel. For \(0<a,a_*<1\), put \(q_a(t)=e^{-2\pi\sqrt t+2\pi ia}\). Then
\begin{equation}\label{da:bilateral-primitive}
 \int_{a_*}^{a}Q_u^{\mathrm{bil}}(t)\dd u
   =\frac{\pi(a-a_*)+
          \Im\{\Li_1(q_a(t))-\Li_1(q_{a_*}(t))\}}{\sqrt t}.
\end{equation}
Indeed \(\partial_a\Li_1(q_a)=2\pi i\Li_0(q_a)\).
The continuous value at \(t=0\) is
\(\pi\{\cot(\pi a_*)-\cot(\pi a)\}\).
Integration of \eqref{da:bilateral-jet} or of its compact connected
version against the shift therefore gives anchored ordinary
polylogarithmic integrals. The required interchanges follow from the
locally uniform bounds already proved for \(Q_a\) and its reflected
partner.
